\begin{claim}{C2.1}{Proved}{Bott--Chern and Aeppli tables as functions of $h$}
With
\[
H^{p,q}_{BC}=
\frac{\ker\partial\cap\ker\bar\partial\subset A^{p,q}}
{\im\partial\bar\partial},\qquad
H^{p,q}_{A}=
\frac{\ker\partial\bar\partial\subset A^{p,q}}
{\im\partial+\im\bar\partial},
\]
the complete tables are
\begin{equation}\label{eq:C-BC-A}
\begin{array}{c|rrrr}
q\backslash p&0&1&2&3\\\hline
3&0&0&1&1\\
2&0&h&2h&1\\
1&0&2&h&0\\
0&1&0&0&0
\end{array}_{BC}
\qquad
\begin{array}{c|rrrr}
q\backslash p&0&1&2&3\\\hline
3&0&0&0&1\\
2&0&h&2&0\\
1&1&2h&h&0\\
0&1&1&0&0.
\end{array}_{A}
\end{equation}
In particular, substitution of the value $h=1$ proved in A3.1 gives
the workbench's numerical tables.  The totals are
$(1,0,2,2,2,2,1)$ and $(1,2,2,2,2,0,1)$, respectively.
\end{claim}
\begin{proof}
We apply the structure theorem for bounded double complexes over a
field, \citet[Theorem 3]{Stelzig2021}, to the actual algebraic double
complex $(A^{\bullet,\bullet}(X),\partial,\bar\partial)$.
The word bounded refers to the bidegrees, not to the dimensions of the
spaces of smooth forms, so that theorem applies here.  It gives a
direct-sum decomposition by isomorphisms of double complexes into
squares and zigzags.  An odd zigzag has one-dimensional total
cohomology.  All odd zigzags in intermediate total degree are therefore
excluded by $b_1=\cdots=b_5=0$.  An odd zigzag of length greater than
one would have at least two vertices of its cohomological total degree.
The extreme diagonals $p+q=0$ and $p+q=6$ contain just one bidegree
each, so the remaining odd zigzags are exactly one dot at $(0,0)$ and
one at $(3,3)$.

An even zigzag of length $2r$ contributes a nonzero $d_r$ to one of
the two Fr\"olicher spectral sequences.  C1.1 and complex conjugation
exclude $r\ge2$ in both sequences.  Thus the remaining summands are
arrows.  The horizontal arrows, with their multiplicities, are
\begin{equation}\label{eq:C-arrows}
\begin{array}{c|c|r}
\text{source}&\text{target}&\text{multiplicity}\\\hline
(0,1)&(1,1)&1\\
(1,1)&(2,1)&h\\
(1,2)&(2,2)&h\\
(2,2)&(3,2)&1.
\end{array}
\end{equation}
The vertical arrows are their reflections in $p=q$, with the same
multiplicities.  The correspondence with spectral-sequence ranks is
\citet[Proposition 6(1)]{Stelzig2021}; alternatively one computes
column cohomology on each horizontal arrow directly.

We now evaluate the displayed definitions, rather than use a formula
for the dimensions.  On an arrow, the source is killed by neither
the outgoing differential nor its total differential, whereas its
target is closed under both.  There is no composite
$\partial\bar\partial$.  Hence the target contributes one dimension
to Bott--Chern cohomology, and the source contributes one dimension
to Aeppli cohomology.  On a square, only the top right corner is closed
under both operators, and it is the image of
$\partial\bar\partial$ from the bottom left corner.  The other
three corners of the square are images of $\partial$ or
$\bar\partial$, while the bottom left corner is not in
$\ker\partial\bar\partial$.  Thus a square contributes to neither
cohomology.  A dot contributes once to each.  These computations
commute with direct sums, including the possibly infinite sum of
squares.  Counting the sources and targets in \eqref{eq:C-arrows}
and its reflection gives \eqref{eq:C-BC-A}.

The program \texttt{code/c\_double\_complex.py} performs a separate
kernel-and-image computation over $\Q$ and a source-and-target
count.  It includes an acyclic square with its signs, checks
$\partial\bar\partial+\bar\partial\partial=0$, and verifies the
two answers for $h=0,1,2,3,4,5$.  Its output is
\texttt{outputs/c\_double\_complex.txt}.  This is exact linear
algebra, not a floating-point check and not a determination of $h$.
\end{proof}
